% =============================================================================
%  CHAPITRE 1 — CADRE CONCEPTUEL ET ÉTAT DE L'ART
% =============================================================================
\entete{Chapitre 1 — Cadre conceptuel et état de l'art}
\chapter{Cadre conceptuel et état de l'art}
\label{chap:cadre}

\introchap

Avant d'exposer la démarche suivie et la solution construite, il importe de
clarifier les notions mobilisées et de situer le travail par rapport à ce qui
existe déjà. Un même mot — « immersion », « visite virtuelle », « assistant
intelligent » — recouvre en effet des réalités très différentes selon les
auteurs et selon les produits commerciaux. Ce chapitre poursuit donc trois
objectifs : définir précisément le vocabulaire employé tout au long du rapport
(section~\ref{sec:conceptuel}), examiner les dispositifs et les travaux
existants afin d'en dégager les apports et les manques
(section~\ref{sec:etatart}), et préciser l'ancrage théorique dans lequel
s'inscrit l'étude (section~\ref{sec:theorique}).

% =============================================================================
\section{Cadre conceptuel}
\label{sec:conceptuel}

\subsection{Patrimoine culturel et valorisation numérique}

Le patrimoine culturel désigne l'ensemble des biens, matériels et immatériels,
qu'une communauté reconnaît comme porteurs de son identité et qu'elle entend
transmettre. La Convention de l'UNESCO (1972) distingue le patrimoine
\emph{matériel} — monuments, ensembles architecturaux, objets — du patrimoine
\emph{immatériel} : savoir-faire, rites, traditions orales. Dans le cas de la
Fondation Jean Félicien Gacha, les deux dimensions coexistent : les objets
exposés ne prennent leur sens que rapportés aux lignées, aux fonctions et aux
usages dont ils procèdent.

La \emph{valorisation numérique} du patrimoine consiste à produire, organiser et
diffuser des représentations numériques de ces biens afin d'en élargir l'accès
et d'en assurer la conservation documentaire. Elle se distingue de la simple
\emph{numérisation}, qui n'est que l'opération technique de conversion d'un
objet ou d'un document en fichier. Numériser, c'est photographier ; valoriser,
c'est rendre intelligible, relier et transmettre. Cette distinction est
essentielle pour notre travail : le problème posé par la Fondation n'était pas
un manque de photographies, mais un manque de mise en relation et de récit.

\subsection{Visite virtuelle, immersion et réalité mixte}

\subsubsection{La visite virtuelle}

Une visite virtuelle est un parcours numérique permettant à un utilisateur
d'explorer un lieu sans s'y rendre physiquement. Trois familles techniques
coexistent :

\begin{itemize}
  \item \textbf{la visite par photographies} : une galerie d'images fixes,
        éventuellement légendées. C'est la forme la plus répandue et la moins
        immersive ; c'est aussi celle que la Fondation pratiquait ;
  \item \textbf{la visite panoramique} : une suite de prises de vue à
        360~degrés, reliées par des points de passage, dans lesquelles
        l'utilisateur pivote librement et se déplace d'un point à l'autre ;
  \item \textbf{la visite en scène tridimensionnelle} : une reconstruction
        complète du lieu, dans laquelle l'utilisateur se déplace librement.
        Elle offre la plus grande liberté mais exige une reconstruction lourde
        et une puissance de calcul importante côté visiteur.
\end{itemize}

\subsubsection{L'immersion}

Slater et Wilbur (1997) proposent de distinguer l'\emph{immersion}, propriété
objective et mesurable du dispositif (champ de vision, latence, fidélité
sensorielle), de la \emph{présence}, état subjectif éprouvé par l'utilisateur
— le sentiment d'« être là ». Un dispositif fortement immersif ne produit pas
mécaniquement un fort sentiment de présence : la cohérence du contenu, la
fluidité de l'interaction et la qualité du récit y contribuent autant que la
technique. Cette distinction guide nos choix : plutôt que de rechercher
l'immersion maximale au prix d'un équipement inaccessible, nous avons cherché le
meilleur sentiment de présence possible sur le matériel dont dispose réellement
le public visé, c'est-à-dire un téléphone ou un ordinateur ordinaire.

\subsubsection{Réalité virtuelle, réalité augmentée, réalité mixte}

Milgram et Kishino (1994) ont formalisé un \emph{continuum} allant du réel
au virtuel :

\begin{itemize}
  \item la \textbf{réalité virtuelle} (VR) remplace intégralement
        l'environnement de l'utilisateur par un environnement de synthèse ;
  \item la \textbf{réalité augmentée} (RA) superpose des éléments de synthèse à
        l'environnement réel, perçu au travers d'une caméra ou d'un verre
        transparent ;
  \item la \textbf{réalité mixte} désigne les situations où le réel et le
        virtuel interagissent, l'objet virtuel étant occulté par les obstacles
        réels et posé sur les surfaces détectées.
\end{itemize}

Notre travail se situe volontairement du côté de la réalité augmentée
« sans équipement » : l'objet est posé dans l'environnement de l'utilisateur au
moyen de la caméra de son téléphone, sans casque ni accessoire. Ce choix est
dicté par le public : exiger un casque reviendrait à recréer, sous une autre
forme, la barrière d'accès que le projet entend supprimer.

\subsection{Modélisation tridimensionnelle et acquisition du réel}

Trois voies permettent d'obtenir le modèle tridimensionnel d'un objet
patrimonial.

\begin{itemize}
  \item \textbf{La modélisation manuelle}, réalisée par un infographiste. Elle
        offre un contrôle total sur le résultat mais coûte cher et ne garantit
        pas la fidélité aux dimensions réelles.
  \item \textbf{La photogrammétrie}, qui reconstruit la géométrie d'un objet à
        partir d'un ensemble de photographies prises sous des angles différents.
        La chaîne classique — détection de points caractéristiques, estimation
        des positions de caméra, nuage de points dense, maillage, texture — est
        implantée par des logiciels tels que Meshroom ou COLMAP. Elle est peu
        coûteuse en matériel mais exigeante en méthode de prise de vue et en
        puissance de calcul.
  \item \textbf{La télémétrie par laser} (LiDAR), qui mesure directement les
        distances. Depuis l'intégration d'un capteur LiDAR dans les téléphones
        haut de gamme, cette technique est devenue accessible et fournit une
        géométrie correctement dimensionnée en quelques minutes.
\end{itemize}

Le format d'échange retenu par l'industrie du web est le glTF, et sa variante
binaire GLB, décrite par le consortium Khronos comme le « JPEG de la 3D ». Le
monde Apple impose en outre un format propre, l'USDZ, pour son dispositif de
prévisualisation. Cette dualité de formats n'est pas un détail : elle conditionne
la couverture réelle des appareils, comme nous le verrons au chapitre~3.

\subsection{Progiciel de gestion intégré et architecture multi-organisations}

Un \emph{progiciel de gestion intégré} (ERP) est un système unique regroupant
les processus de gestion d'une organisation autour d'une base de données
partagée. Appliqué à une institution patrimoniale, il couvre l'inventaire des
collections, la gestion des espaces, la billetterie, la boutique, les événements
et la relation avec le public.

Le modèle du \emph{logiciel en tant que service} (SaaS) consiste à faire
fonctionner une seule instance de l'application pour plusieurs organisations
clientes. On parle alors d'architecture \emph{multi-organisations}
(\emph{multi-tenant}). Chong et Carraro (2006) en distinguent trois
implantations : une base de données par organisation, un schéma par organisation,
ou une base commune avec cloisonnement par ligne. La troisième est la plus
économique et la plus délicate : elle exige que la base elle-même, et non le
code applicatif, garantisse qu'une organisation ne puisse jamais lire les
données d'une autre. C'est le rôle de la \emph{sécurité au niveau de la ligne}
(RLS), mécanisme du SGBD qui filtre les lignes visibles selon l'identité de
l'utilisateur, indépendamment de la requête envoyée.

\subsection{Intelligence artificielle générative et génération augmentée par la recherche}

\subsubsection{Grands modèles de langue}

Un grand modèle de langue est un réseau de neurones entraîné à prédire la suite
d'un texte. L'architecture dominante, le \emph{Transformer}, a été introduite par
Vaswani et al.\ (2017). Ces modèles produisent un texte fluide et
grammaticalement correct, mais ils ne disposent d'aucun mécanisme interne de
vérification factuelle : ils peuvent énoncer avec assurance des faits inexacts,
phénomène désigné dans la littérature sous le terme d'\emph{hallucination}
(Ji et al., 2023). Dans un contexte de médiation culturelle, où l'exactitude
prime, ce défaut est rédhibitoire s'il n'est pas corrigé.

\subsubsection{Plongements lexicaux et recherche sémantique}

Un \emph{plongement} (\emph{embedding}) est la représentation d'un texte sous
forme de vecteur numérique, construite de telle manière que deux textes de sens
voisin correspondent à deux vecteurs proches. La proximité se mesure
habituellement par la similarité cosinus. Cette représentation permet une
recherche par le sens et non par les mots : une requête portant sur un
« masque de danse » peut ainsi faire remonter une notice mentionnant un
« masque cérémoniel », alors qu'aucun mot n'est commun aux deux formulations.

Le stockage et l'interrogation de ces vecteurs demandent un index adapté. La
structure HNSW, proposée par Malkov et Yashunin (2018), permet une recherche
approchée des plus proches voisins en temps quasi logarithmique, ce qui rend la
recherche sémantique praticable sur des dizaines de milliers de documents.

\subsubsection{La génération augmentée par la recherche}

La génération augmentée par la recherche (RAG), formalisée par Lewis et al.\
(2020), consiste à ne pas demander au modèle de répondre de mémoire, mais à lui
fournir, au moment de la question, les documents pertinents extraits d'une base
maîtrisée, en lui imposant de fonder sa réponse sur eux seuls. Le mécanisme se
décompose en cinq temps :

\begin{enumerate}
  \item \textbf{l'ingestion} : les documents de référence sont collectés,
        découpés en fragments de taille homogène et stockés ;
  \item \textbf{l'indexation} : chaque fragment est converti en vecteur par un
        modèle de plongement, puis enregistré dans un index vectoriel ;
  \item \textbf{la recherche} : la question de l'utilisateur est convertie en
        vecteur et comparée à l'index, qui renvoie les fragments les plus
        proches ;
  \item \textbf{l'augmentation} : les fragments retenus sont insérés dans
        l'instruction envoyée au modèle de langue, assortis de la consigne de ne
        rien affirmer qui n'y figure ;
  \item \textbf{la génération} : le modèle rédige la réponse et cite les
        sources ; s'il ne trouve pas l'information, il doit le dire.
\end{enumerate}

L'intérêt de ce dispositif pour une institution patrimoniale est direct : il
permet d'obtenir la fluidité d'un guide et la fiabilité d'un catalogue, et il
rend la mise à jour immédiate — modifier une notice suffit à modifier les
réponses, sans réentraîner quoi que ce soit.

\subsection{DevOps, conteneurisation et infrastructure décrite par le code}

\subsubsection{DevOps}

Le terme DevOps désigne un ensemble de pratiques visant à raccourcir le cycle
qui va de l'écriture d'une modification à sa mise à disposition des
utilisateurs, en supprimant la frontière entre l'équipe qui développe et celle
qui exploite. Kim et al.\ (2016) en résument le principe par « trois voies » :
fluidifier le flux du développement vers la production, raccourcir les boucles
de retour, et instaurer une culture d'expérimentation et d'apprentissage
continus. Deux pratiques en constituent le cœur opérationnel :

\begin{itemize}
  \item \textbf{l'intégration continue} : chaque modification déposée déclenche
        automatiquement la compilation et une série de vérifications ; une
        anomalie est signalée en quelques minutes, et non le jour de la mise en
        ligne ;
  \item \textbf{le déploiement continu} : une modification validée est publiée
        automatiquement, selon une procédure identique à chaque fois, ce qui
        supprime les erreurs de manipulation.
\end{itemize}

\subsubsection{Conteneurisation}

Un \emph{conteneur} encapsule une application avec l'intégralité de ce dont elle
a besoin pour fonctionner — bibliothèques, fichiers de configuration, serveur —
de sorte qu'elle se comporte de manière identique sur le poste du développeur et
sur le serveur de production. Docker en est l'implantation de référence. Le
conteneur résout le problème classiquement résumé par la formule « chez moi, ça
marche » : ce n'est plus le code seul qui est transporté, mais son environnement
complet.

Lorsque plusieurs conteneurs doivent être lancés, surveillés, remplacés et
répartis sur plusieurs machines, un \emph{orchestrateur} devient nécessaire. Les
deux services étudiés dans ce travail sont ECS, orchestrateur propre à AWS, plus
simple à mettre en œuvre, et EKS, distribution gérée de Kubernetes, plus riche
mais plus coûteuse en exploitation.

\subsubsection{Infrastructure décrite par le code}

L'\emph{infrastructure décrite par le code} (IaC) consiste à définir dans des
fichiers versionnés l'ensemble des ressources à créer chez le fournisseur de
nuage — stockage, réseau de diffusion, certificats, noms de domaine, droits —
puis à laisser un outil les créer et les maintenir conformes. Terraform en est
l'outil le plus répandu. L'apport est double : l'infrastructure devient
reproductible à l'identique, et toute modification est tracée et relisible comme
du code.

\subsection{FinOps : le pilotage économique du nuage}

Le nuage transforme une dépense d'investissement en dépense de fonctionnement :
on ne paie plus un serveur, on paie ce que l'on consomme. Cette souplesse a une
contrepartie — la dépense devient variable, diffuse et facilement invisible. La
discipline dite \emph{FinOps} vise à rendre ce coût visible, à l'attribuer à
ceux qui le provoquent et à l'optimiser en continu. Elle repose sur trois
temps : \emph{informer} (mesurer et répartir la dépense), \emph{optimiser}
(supprimer le gaspillage, dimensionner au juste besoin) et \emph{opérer}
(inscrire ces contrôles dans le fonctionnement courant). Pour une institution à
budget contraint comme la Fondation, cette discipline n'est pas un raffinement :
elle conditionne la possibilité même de maintenir le service dans la durée.

% =============================================================================
\section{État de l'art}
\label{sec:etatart}

\subsection{Les dispositifs de visite virtuelle existants}

Le marché de la visite virtuelle s'est structuré autour de quelques familles de
solutions, dont les caractéristiques sont résumées au tableau~\ref{tab:comparvv}.

\begin{table}[htbp]
\centering
\caption{Comparaison des principales familles de solutions de visite virtuelle}
\label{tab:comparvv}
\small
\begin{tabularx}{\textwidth}{|L{3.1cm}|X|L{3.0cm}|L{2.6cm}|}
\hline
\rowcolor{keycetete}
\thead{Solution} & \thead{Principe et apports} & \thead{Limites relevées} & \thead{Modèle économique} \\ \hline
Google Arts \& Culture &
Numérisation à très haute définition et parcours en vue immersive ; audience
mondiale considérable &
Sélection à l'initiative de l'opérateur ; l'institution ne maîtrise ni sa
présentation ni ses données ; aucun outil de gestion &
Gratuit, sur invitation \\ \hline
Matterport &
Capture tridimensionnelle d'un lieu par caméra dédiée ; rendu de grande qualité,
« maison de poupée » navigable &
Matériel spécifique coûteux ; hébergement des données chez l'éditeur ; pas
d'assistance conversationnelle ni de gestion documentaire &
Abonnement mensuel par espace \\ \hline
Kuula, Pano2VR &
Assemblage de panoramas à 360~degrés avec points d'intérêt ; prise en main
rapide et coût faible &
Pas de base documentaire, pas de réalité augmentée, pas de gestion de
l'institution ; parcours isolé du reste du site &
Abonnement, palier gratuit \\ \hline
Développement sur mesure &
Adapté au besoin exact de l'institution &
Coût de conception élevé, non réutilisable pour une autre structure, maintenance
à la charge du commanditaire &
Prestation ponctuelle \\ \hline
\end{tabularx}
\source{Élaboré par nos soins à partir de la documentation publique des éditeurs
(consultée en juin 2026).}
\end{table}

Il ressort de cette comparaison une constante : les solutions existantes traitent
la visite virtuelle comme un \emph{produit isolé}, détaché de la gestion de
l'institution et de sa documentation. Aucune ne relie le parcours immersif à un
inventaire, à une billetterie, à une base de connaissances interrogeable ou à
d'autres institutions. C'est précisément cette articulation qui fait défaut et
que notre travail se propose de construire.

\subsection{Les collections ouvertes et leurs interfaces de programmation}

Plusieurs grands musées ont ouvert leurs catalogues au moyen d'interfaces de
programmation publiques. Ces sources constituent la matière première du
rapprochement documentaire recherché. Le tableau~\ref{tab:collections} présente
celles qui ont été examinées.

\begin{table}[htbp]
\centering
\caption{Collections ouvertes examinées et conditions d'accès}
\label{tab:collections}
\small
\begin{tabularx}{\textwidth}{|L{4.0cm}|C{2.2cm}|X|}
\hline
\rowcolor{keycetete}
\thead{Institution ou service} & \thead{Clé requise} & \thead{Observation} \\ \hline
The Metropolitan Museum of Art & Non & Interface publique, bonne couverture africaine ; source principale retenue \\ \hline
Art Institute of Chicago (ARTIC) & Non & Interface documentée et stable ; retenue \\ \hline
Cleveland Museum of Art & Non & Notices détaillées, numéros d'inventaire fournis ; retenue \\ \hline
Victoria and Albert Museum & Non & Riche en arts appliqués ; retenue \\ \hline
Wikidata & Non & Données structurées reliant objets, cultures et lieux ; retenue comme source de liaison \\ \hline
Europeana & \textbf{Oui} & Couverture européenne considérable ; écartée faute de clé d'accès \\ \hline
Rijksmuseum & \textbf{Oui} & Écartée pour la même raison \\ \hline
Smithsonian Institution & \textbf{Oui} & Écartée pour la même raison \\ \hline
Harvard Art Museums & \textbf{Oui} & Écartée pour la même raison \\ \hline
\end{tabularx}
\source{Relevé effectué par nos soins en juillet 2026 sur la documentation
technique publiée par chaque institution.}
\end{table}

Ce relevé appelle une remarque de méthode : le choix des sources n'a pas été
dicté par leur richesse, mais par leur accessibilité effective dans les
conditions du stage. Quatre institutions majeures ont dû être écartées faute de
clé d'accès, ce qui constitue une limite explicite de l'étude et non un oubli.

\subsection{Les assistants conversationnels en contexte muséal}

Les premiers agents conversationnels muséaux reposaient sur des arbres de
questions prédéfinies : robustes mais rigides, ils ne répondaient qu'aux
questions prévues. L'arrivée des grands modèles de langue a permis une
conversation naturelle, au prix du risque d'invention évoqué plus haut. La
littérature récente converge vers l'ancrage documentaire comme réponse à ce
risque. Trois exigences sont régulièrement formulées :

\begin{itemize}
  \item \textbf{la traçabilité} : chaque affirmation doit pouvoir être rapportée
        à un document identifiable ;
  \item \textbf{l'aveu d'ignorance} : l'assistant doit pouvoir répondre qu'il ne
        sait pas, ce qui suppose que la consigne l'y autorise explicitement ;
  \item \textbf{le périmètre} : l'assistant doit refuser poliment les questions
        hors sujet plutôt que d'y répondre approximativement.
\end{itemize}

Ces trois exigences ont directement structuré la conception de notre assistant,
décrite au chapitre~3.

\subsection{Positionnement de notre travail}

Au terme de cet examen, trois manques se dégagent, qui définissent l'espace dans
lequel s'inscrit notre contribution.

\begin{enumerate}
  \item \textbf{L'absence d'articulation.} Les dispositifs existants séparent la
        visite immersive, la gestion de l'institution et l'assistance au
        visiteur. Nous proposons de les réunir dans un ensemble unique où un
        objet du parcours immersif est le même objet que celui de l'inventaire,
        de la boutique et de la base de connaissances de l'assistant.

  \item \textbf{L'absence de mise en relation entre institutions.} Aucune
        solution examinée ne cherche, pour une œuvre donnée, ses œuvres
        apparentées ailleurs dans le monde. Nous proposons un dispositif de
        rapprochement documentaire assorti d'une validation humaine.

  \item \textbf{L'absence de modèle reproductible.} Les solutions sur mesure ne
        se transposent pas d'une institution à l'autre. Nous proposons une
        architecture multi-organisations dans laquelle l'accueil d'une nouvelle
        structure ne demande aucun développement.
\end{enumerate}

% =============================================================================
\section{Ancrage théorique}
\label{sec:theorique}

Bien que le guide méthodologique autorise les travaux de licence à s'en tenir au
cadre conceptuel, il nous a paru utile de préciser les deux cadres théoriques
qui orientent l'interprétation de nos résultats.

Le premier est la \textbf{recherche par la conception} (\emph{design science
research}), formalisée par Hevner et al.\ (2004). Elle établit que la production
d'un artefact — modèle, méthode ou système — constitue une contribution
scientifique légitime, à condition que l'artefact réponde à un problème réel,
qu'il soit évalué avec rigueur et que les enseignements tirés soient
généralisables. Ce cadre justifie la nature même de notre démarche : nous ne
nous contentons pas de décrire une situation, nous construisons un dispositif et
nous l'évaluons.

Le second est le \textbf{modèle d'acceptation des technologies} (TAM), proposé
par Davis (1989), selon lequel l'adoption d'un outil dépend principalement de
son utilité perçue et de sa facilité d'utilisation perçue. Ce modèle éclaire
plusieurs de nos arbitrages : le refus du casque de réalité virtuelle, le choix
d'un accès par simple navigateur, ou encore la mise à disposition d'un code QR
permettant de passer de l'ordinateur au téléphone sans installation, sont autant
de décisions visant à maximiser la facilité d'utilisation perçue auprès d'un
public non technicien.

\conclchap

Ce chapitre a précisé le vocabulaire du rapport et situé le travail dans son
paysage. Il en ressort que les briques techniques nécessaires — panoramas à
360~degrés, modélisation tridimensionnelle, réalité augmentée, plongements
lexicaux, génération augmentée par la recherche, cloisonnement par la base de
données — sont individuellement disponibles et éprouvées. Ce qui manque n'est
donc pas la technologie, mais leur \emph{assemblage} au service d'une
institution patrimoniale à moyens limités. C'est cet assemblage qui constitue
l'objet de notre travail. Le chapitre suivant expose la démarche méthodologique
adoptée pour le concevoir, le construire et l'évaluer.
